\subsection{An explicit obstruction to the proposed Hooley bridge}

\paragraph{Finding A06 (the universal primitive-root theta condition is false).}
The condition $T(C)$ of A03 does not hold for every genus-two curve in
characteristic two, even after arbitrary finite extension of its field of
definition. Take the smooth proper model of
\[
 \begin{split}
 k&=\mathbb F_2[t]/(M(t)),\quad
 M(Z)=Z^{38}+Z^{14}+Z^{10}+Z^9+Z^8+Z^5+Z^2+Z+1,\\
 d&=t^3+1,\qquad
 C:\ y^2+y=A(u)^3,\qquad
 A(u)=(u^2+t^2u+t)/d.
 \end{split} \tag{A06.1}
\]
For its locally exact differential theta divisor $\Theta$, there is no
odd prime $\ell$ with $\operatorname{ord}_\ell(2)=\ell-1$ and
$\Theta\cap J_C[\ell](\bar k)\ne\{0\}$. This holds geometrically;
number-field GRH cannot change it.

\begin{proof}
The curve and scalar Frobenius are those of \cite[B.23]{A1run}.
We recall the algebraic data used here and give the new small
primitive-root obstruction in full. Rabin's irreducibility test gives
$Z^{2^{38}}=Z$ modulo $M$ and
$\gcd(M,Z^{2^2}-Z)=\gcd(M,Z^{2^{19}}-Z)=1$.
Thus $k$ is the stated field and $d\ne0$. Reduction of the
Artin--Schreier polynomial by even powers leaves leading odd term
$(t^2/d^3)u^5$. There is a unique point $P$ at infinity, the genus
is two, and $K_C\sim2P$. The one-pole Artin--Schreier curve has
two-rank zero (equivalently, Cartier on the basis $du,u\,du$ has
stable rank zero). Hence $J_C[2](\bar k)=0$ and the unique geometric
theta characteristic is $\mathcal O(P)$. In characteristic two
the bundle of locally exact differentials is this theta characteristic,
so its theta divisor is $\Theta=C-P$. With relative Frobenius this
statement is on $C^{(1)}$; the same equations with coefficients squared
and the same scalar Frobenius give exactly the argument below.
In particular it is symmetric
and contains zero.

Put $B(u)=(tu^2+u+t^2)/d$. There are two maps to
$E:y^2+y=x^3$, given by $(A(u),y)$ and
$(B(u),y+(u^3+1)/d)$. Their equations follow from
\[
 (u^2+t^2u+t)^3+(tu^2+u+t^2)^3
 =d(u^3+1)^2+d^2(u^3+1).
\]
Their joint map is birational onto its image, since $u=B+tA$.
The image, whose normalization has genus two, cannot lie in an
elliptic translate in $E^2$; thus $J_C\to E^2$ is an isogeny.
Both maps take $P$ to zero. The duplication formula on $E$ is
$[2](x,y)=(x^4,y^4+1)$, so its fourth-power Frobenius is $[-2]$.
It follows that the $k$-Frobenius of $J_C$ is $[-2^{19}]$.
Indeed composing the difference with the isogeny gives zero,
and a homomorphism from a connected abelian variety into a finite
kernel is zero. These geometric identities and the irreducibility
calculation were separately recomputed in this run.

A nonzero difference in $J_C$ has at most two ordered representations
as a difference of points on $\Theta$. For if
$R_1-S_1=R_2-S_2\ne0$, then
$R_1+S_2\sim R_2+S_1$. Either the effective divisors coincide,
which gives the same representation, or they belong to the unique
degree-two hyperelliptic pencil. The latter determines the sole
alternative representation $(\iota S_1,\iota R_1)$.

Suppose a primitive-root prime $\ell$ and a nonzero
$L\in\Theta[\ell]$ existed. Since the theta divisor is defined over
$k$ and symmetric, it contains $hL$ for
$h\in H=\langle2^{19}\rangle\subset\mathbb F_\ell^\times$.
Write $e=\gcd(\ell-1,19)$ and $m=(\ell-1)/e$.
It contains these $m$ distinct points and zero. Counting their
ordered nonzero differences in the cyclic subgroup gives
\[
 m(m+1)\le2(\ell-1). \tag{A06.2}
\]
If $e=1$, this forces $\ell\le2$, impossible. If $e=19$,
it forces $m\le37$, hence $\ell\le704$ and
$\ell\equiv1\pmod{19}$. Exhaustive trial division of the eighteen
integers $19m+1$ with $m=2,4,\ldots,36$ gives exactly
\[
\begin{array}{c|rrrrrr}
 \ell&191&229&419&457&571&647\\ \hline
 \operatorname{ord}_\ell(2)&95&76&418&76&114&323.
\end{array}
\]
Only $419$ remains. Here $g=2^{19}\bmod419=119$, and $H$ contains
\[
 (g^0,g^{11})=(1,418),\quad
 (g^{19},g^{20})=(71,69),\quad
 (g^9,g^8)=(350,348).
\]
These are three distinct ordered pairs with the same nonzero
difference $2$ modulo $419$, contradicting the two-representation
bound. The included standard-library certificate
\texttt{check\_theta\_obstruction.py} verifies every finite calculation.
\end{proof}

\paragraph{Scope of the refutation.}
The stronger source statement \cite[B.23]{A1run} excludes every
nontrivial odd cyclic etale cover with zero primitive two-rank on this
curve. Its exhaustive conductor reduction and its complete polynomial
B\'ezout certificates at $191$ and $229$ were rerun successfully;
the verifier, data and logs are included. The shorter A06 proof does
not depend on those large certificates. In particular the weaker
arithmetic condition $\langle2,-1\rangle=\mathbb F_\ell^\times$
cannot repair the universal geometric assertion on this example,
by that checked B.23 result.

This curve already has two-rank zero, so ordinary Kummer theory
separates its points. It is not a counterexample to A1, nor a
refutation of a strategy which first removes all zero-rank curves
and then seeks a different theorem for the remaining curves.
The positive-two-rank theta-incidence problem has not been solved
by this argument. Its unresolved status does not affect A05.
